\documentclass[conference,10pt]{IEEEtran}

\usepackage[T1]{fontenc}
\usepackage[utf8]{inputenc}
\usepackage[spanish,es-nodecimaldot]{babel}
\usepackage{amsmath,amssymb}
\usepackage{booktabs}
\usepackage{xcolor}
\usepackage{float}
\usepackage{graphicx}
\usepackage[american]{circuitikz}
\usetikzlibrary{babel}

\definecolor{placeholder}{RGB}{110,110,110}
\newcommand{\completar}[1]{%
  \ifmmode\text{\color{placeholder}\itshape[Completar: #1]}%
  \else\textcolor{placeholder}{\textit{[Completar: #1]}}\fi}
\newcommand{\espaciolargo}{\par\vspace{1.5cm}}
\newcommand{\espaciocalculos}{\par\vspace{3.2cm}}

\title{Preinforme: Circuitos Resistivos - Parte 2\\
Práctica de Laboratorio No. 3}

\author{
  \IEEEauthorblockN{Maria Paula Mejía Vásquez, Cristian Alejandro Guarin Marin}
  \IEEEauthorblockA{
    Ingeniería Electrónica y de Telecomunicaciones\\
    Circuitos Eléctricos I -- Universidad de Antioquia\\
    Periodo 2026-2\\
    paula.mejia3@udea.edu.co, cristian.guarinm@udea.edu.co
  }
}

\begin{document}
\maketitle

\begin{abstract}
Este preinforme prepara la práctica de circuitos resistivos (Parte 2), enfocada en la verificación teórica de las Leyes de Kirchhoff de Corriente (LCK) y Voltaje (LVK). Se seleccionan los componentes requeridos y se desarrollan los cálculos teóricos empleando el análisis nodal y el análisis de mallas para un circuito de dos lazos. Los resultados obtenidos permitirán validar matemáticamente el cumplimiento de las leyes de conservación de energía y carga, sirviendo como base de referencia para el posterior montaje físico y la simulación en LTspice.
\end{abstract}

\begin{IEEEkeywords}
Leyes de Kirchhoff, análisis nodal, análisis de mallas, circuitos resistivos, simulación.
\end{IEEEkeywords}

\section{Introducción}
El análisis sistemático de circuitos eléctricos requiere el dominio de herramientas matemáticas que permitan determinar voltajes y corrientes en cualquier punto de una red. Las Leyes de Kirchhoff, junto con la Ley de Ohm, constituyen el pilar de los métodos de análisis nodal (basado en la LCK) y de mallas (basado en la LVK). 

En esta práctica se analiza teóricamente una configuración resistiva de dos mallas y cuatro nodos. A través del cálculo previo de voltajes nodales y corrientes de lazo, es posible deducir el comportamiento detallado de cada componente, comprobando que la sumatoria de caídas de tensión en trayectorias cerradas es nula y que el flujo neto de corriente en las uniones es cero.

\section{Objetivos}
\subsection{Objetivo general}
Analizar de forma matemática un circuito resistivo aplicando las leyes de Kirchhoff de corriente (LCK) y de voltaje (LVK) a través de los métodos de mallas y nodos, estableciendo los valores teóricos previos a la verificación experimental.

\subsection{Objetivos específicos}
\begin{itemize}
  \item Seleccionar tres resistores comerciales diferentes en el rango de $680~\Omega$ a $2.2~\mathrm{k}\Omega$ a $1/4~\mathrm{W}$ para el montaje.
  \item Determinar matemáticamente el voltaje en todos los nodos del circuito empleando análisis nodal con referencia en tierra.
  \item Hallar las corrientes de malla del circuito planteando las ecuaciones del análisis de lazos.
  \item Verificar teóricamente el cumplimiento de la LVK en cada malla y de la LCK en cada nodo a partir de los resultados calculados.
\end{itemize}

\section{Componentes}
\subsection{Materiales}
Para el desarrollo de la práctica se emplearán los siguientes elementos:
\begin{itemize}
  \item Protoboard y cables de conexión.
  \item Multímetro.
  \item Fuente de voltaje dual.
  \item Resistores comerciales.
\end{itemize}

\subsection{Valores de los componentes}
De acuerdo con la guía, se seleccionan tres resistores en el rango solicitado y se establecen los valores de las fuentes.

\begin{table}[H]
  \caption{Valores nominales de los componentes para el laboratorio}
  \label{tab:componentes}
  \centering
  \resizebox{\columnwidth}{!}{%
  \begin{tabular}{ccc}
    \toprule
    Componente & Valor Nominal & Especificación \\
    \midrule
    Resistor $R_1$ & $680~\Omega$ & $1/4~\mathrm{W}$ \\
    Resistor $R_2$ & $1000~\Omega$ & $1/4~\mathrm{W}$ \\
    Resistor $R_3$ & $2000~\Omega$ & $1/4~\mathrm{W}$ \\
    Fuente $v_1$ & $8~\mathrm{V}$ & Corriente directa \\
    Fuente $v_2$ & $6~\mathrm{V}$ & Corriente directa \\
    \bottomrule
  \end{tabular}%
  }
\end{table}

\section{Cálculos}

\begin{figure}[H]
  \centering
  \begin{circuitikz}[american, scale=0.8, transform shape]
    % Nodos inferiores (d)
    \draw (0,0) -- (6,0) node[pos=0.5, below] {$d$};
    % Lazo izquierdo
    \draw (0,0) to[battery1, invert, l=$v_1$, v_<=$8\mathrm{V}$] (0,3) node[above] {$a$};
    \draw (0,3) to[R, l=$R_1$] (3,3) node[above] {$b$};
    \draw (3,3) to[R, l=$R_2$] (3,0);
    % Lazo derecho
    \draw (3,3) to[battery1, l=$v_2$, v_>=$6\mathrm{V}$] (6,3) node[above] {$c$};
    \draw (6,3) to[R, l=$R_3$] (6,0);

    % Corrientes de malla (horario y sin superponerse)
    \draw[->, red, thick] (1.0, 1.5) arc (180:-90:0.5);
    \node[red] at (1.5, 1.5) {$i_1$};

    \draw[->, red, thick] (4.0, 1.5) arc (180:-90:0.5);
    \node[red] at (4.5, 1.5) {$i_2$};
  \end{circuitikz}
  \caption{Circuito para la realización de la práctica.}
  \label{fig:circuito-practica}
\end{figure}

\subsection{Tabla de voltajes de nodo y corrientes de malla (numeral a)}
En la siguiente tabla se resumen los voltajes de nodo y las corrientes de malla teóricas, halladas mediante el análisis nodal y el análisis de mallas desarrollados en los numerales siguientes.

\begin{table}[H]
  \caption{Voltajes de nodo y corrientes de malla - preinforme}
  \label{tab:resultados}
  \centering
  \resizebox{\columnwidth}{!}{%
  \begin{tabular}{cccccc}
    \toprule
    \multicolumn{4}{c}{Voltajes de nodo (V)} & \multicolumn{2}{c}{Corrientes de malla (mA)} \\
    \cmidrule(r){1-4} \cmidrule(l){5-6}
    $v_a$ & $v_b$ & $v_c$ & $v_d$ & $i_1$ & $i_2$ \\
    \midrule
    $8.000$ & $4.970$ & $-1.030$ & $0$ & $4.455$ & $-0.5148$ \\
    \bottomrule
  \end{tabular}%
  }
\end{table}

\subsection{Análisis nodal (numeral b)}
Se ubica la referencia (tierra) en el nodo $d$, de modo que $v_d=0~\mathrm{V}$. Como la fuente $v_1$ conecta directamente el nodo $a$ con la referencia,
\begin{equation}
  v_a = v_1 = 8~\mathrm{V}.
\end{equation}

La fuente $v_2$ conecta los nodos $b$ y $c$ sin una resistencia en serie entre ellos, por lo que estos dos nodos forman un supernodo, cuya restricción de tensión es
\begin{equation}
  v_b - v_c = v_2 = 6~\mathrm{V}. \label{eq:restriccion-supernodo}
\end{equation}

Se definen las corrientes de rama $i_{R1}$ (del nodo $a$ hacia el supernodo, a través de $R_1$), $i_{R2}$ (del supernodo hacia $d$, a través de $R_2$) e $i_{R3}$ (del supernodo hacia $d$, a través de $R_3$). Aplicando la LCK al supernodo (la corriente que entra es igual a la suma de las corrientes que salen),
\begin{equation}
  i_{R1} = i_{R2} + i_{R3}. \label{eq:lck-supernodo}
\end{equation}

Empleando la Ley de Ohm en cada rama, $i_{R1}=\dfrac{v_a-v_b}{R_1}$, $i_{R2}=\dfrac{v_b}{R_2}$ e $i_{R3}=\dfrac{v_c}{R_3}$, y sustituyendo en~\eqref{eq:lck-supernodo},
\begin{equation}
  \frac{v_a-v_b}{R_1} = \frac{v_b}{R_2} + \frac{v_c}{R_3}. \label{eq:supernodo}
\end{equation}

Sustituyendo $v_a=8~\mathrm{V}$ y los valores de los resistores en~\eqref{eq:supernodo},
\begin{equation}
  \frac{8-v_b}{680} = \frac{v_b}{1000} + \frac{v_c}{2000}.
\end{equation}
Multiplicando ambos lados por $2000$ y agrupando términos,
\begin{equation}
  \frac{2000}{680}(8-v_b) = 2v_b + v_c \;\Rightarrow\; 4.941\,v_b + v_c = 23.529.
\end{equation}

Junto con la restricción del supernodo~\eqref{eq:restriccion-supernodo}, se obtiene el sistema de ecuaciones en forma matricial:
\begin{equation*}
  \begin{bmatrix}
    4.941 & 1 \\
    1 & -1
  \end{bmatrix}
  \begin{bmatrix}
    v_b \\
    v_c
  \end{bmatrix}
  =
  \begin{bmatrix}
    23.529 \\
    6
  \end{bmatrix}
\end{equation*}

Resolviendo el sistema matricial, se obtienen los voltajes de nodo:
\begin{equation}
  v_b = 4.970~\mathrm{V}, \qquad v_c = v_b - 6 = -1.030~\mathrm{V}.
\end{equation}

\noindent\textbf{Resultados:}
\begin{align*}
  v_a &= 8~\mathrm{V}, & v_b &= 4.970~\mathrm{V},\\
  v_c &= -1.030~\mathrm{V}, & v_d &= 0~\mathrm{V}.
\end{align*}

\subsection{Análisis de mallas (numeral c)}
Se determinan las corrientes de malla del circuito ($i_1, i_2$) planteando las ecuaciones de lazo correspondientes.

\textbf{Malla 1:}
\begin{align*}
    -v_1 + R_1 i_1 + R_2 (i_1 - i_2) &= 0 \\
    (R_1 + R_2) i_1 - R_2 i_2 &= v_1 \\
    (680 + 1000) i_1 - 1000 i_2 &= 8 \\
    1680 i_1 - 1000 i_2 &= 8 \quad \text{(Eq. 1)}
\end{align*}

\textbf{Malla 2:}
\begin{align*}
    R_2 (i_2 - i_1) + v_2 + R_3 i_2 &= 0 \\
    -R_2 i_1 + (R_2 + R_3) i_2 &= -v_2 \\
    -1000 i_1 + (1000 + 2000) i_2 &= -6 \\
    -1000 i_1 + 3000 i_2 &= -6 \quad \text{(Eq. 2)}
\end{align*}

Expresando el sistema de ecuaciones en forma matricial:
\begin{equation*}
    \begin{bmatrix}
        1680 & -1000 \\
        -1000 & 3000
    \end{bmatrix}
    \begin{bmatrix}
        i_1 \\
        i_2
    \end{bmatrix}
    =
    \begin{bmatrix}
        8 \\
        -6
    \end{bmatrix}
\end{equation*}

Resolviendo el sistema matricial, se obtienen las corrientes de malla:
\begin{align*}
    i_1 &= 0.004455~\mathrm{A} = 4.455~\mathrm{mA} \\
    i_2 &= -0.0005148~\mathrm{A} = -0.5148~\mathrm{mA}
\end{align*}

\subsection{Verificación de la Ley de Voltajes de Kirchhoff (numeral d)}
Empleando los voltajes de nodo hallados en el numeral anterior, se determinan los voltajes en cada componente:
\begin{align}
  v_1 &= v_a - v_d = 8 - 0 = 8~\mathrm{V},\\
  v_{R1} &= v_a - v_b = 8 - 4.970 = 3.030~\mathrm{V},\\
  v_{R2} &= v_b - v_d = 4.970 - 0 = 4.970~\mathrm{V},\\
  v_2 &= v_b - v_c = 4.970-(-1.030) = 6~\mathrm{V},\\
  v_{R3} &= v_c - v_d = -1.030 - 0 = -1.030~\mathrm{V}.
\end{align}

\textbf{Malla 1:} recorriendo el lazo $a$--$b$--$d$--$a$ en el sentido de $i_1$,
\begin{equation}
  \sum v_{\text{malla}1} = -v_1 + v_{R1} + v_{R2} = -8 + 3.030 + 4.970 = 0.
\end{equation}

\textbf{Malla 2:} recorriendo el lazo $b$--$c$--$d$--$b$ en el sentido de $i_2$,
\begin{equation}
  \sum v_{\text{malla}2} = -v_{R2} + v_2 + v_{R3} = -4.970 + 6 - 1.030 = 0.
\end{equation}

Se comprueba que la suma algebraica de las caídas de tensión en cada malla es igual a cero, verificando la Ley de Voltajes de Kirchhoff.

\subsection{Verificación de la Ley de Corrientes de Kirchhoff (numeral e)}
Empleando las corrientes de malla halladas, se deducen las corrientes de cada rama y se verifica que la sumatoria en cada nodo principal es cero ($\sum i_{in} - \sum i_{out} = 0$).

Previamente, hallamos la corriente en la rama compartida ($R_2$):
\begin{equation*}
    i_{R2} = i_1 - i_2 = 4.455~\mathrm{mA} - (-0.5148~\mathrm{mA}) = 4.9698~\mathrm{mA}
\end{equation*}

\textbf{Nodo a (LCK):}
\begin{align*}
    \sum i_a &= i_{v1} - i_{R1} = i_1 - i_1 \\
    \sum i_a &= 4.455 - 4.455 = 0
\end{align*}

\textbf{Nodo b (LCK):}
\begin{align*}
    \sum i_b &= i_1 - i_{R2} - i_2 = 0 \\
    \sum i_b &= 4.455 - (4.455 - (-0.5148)) - (-0.5148) \\
    \sum i_b &= 4.455 - 4.9698 + 0.5148 = 0
\end{align*}

\textbf{Nodo c (LCK):}
\begin{align*}
    \sum i_c &= i_2 - i_{R3} = i_2 - i_2 \\
    \sum i_c &= -0.5148 - (-0.5148) = 0
\end{align*}

\textbf{Nodo d (LCK):}
\begin{align*}
    \sum i_d &= i_2 + i_{R2} - i_1 = 0 \\
    \sum i_d &= (-0.5148) + 4.9698 - 4.455 = 0
\end{align*}

\section{Referencias}
\begin{thebibliography}{00}
  \bibitem{guia}
  Universidad de Antioquia, ``Guía de Práctica de Laboratorio No. 3: Circuitos Resistivos - Parte 2,'' Departamento de Ingeniería Electrónica y de Telecomunicaciones, Medellín, Colombia, 2026.
  \bibitem{texto}
  C. K. Alexander y M. N. O. Sadiku, \emph{Fundamentos de circuitos eléctricos}, 5.ª ed. México D.F., México: McGraw-Hill Interamericana, 2013.
\end{thebibliography}

\end{document}
